\documentclass[11pt]{article}

\usepackage[margin=1in]{geometry}
\usepackage{microtype}
\usepackage{parskip}
\usepackage[hidelinks]{hyperref}
\usepackage{titlesec}
\usepackage{enumitem}
\usepackage{booktabs}
\usepackage{amsmath}
\usepackage{amssymb}

\setlist{nosep, leftmargin=1.5em, topsep=2pt}
\titleformat{\section}{\large\bfseries}{\thesection.}{0.6em}{}
\titleformat{\subsection}{\normalsize\bfseries}{\thesubsection}{0.6em}{}
\titlespacing*{\section}{0pt}{1.6em}{0.5em}
\titlespacing*{\subsection}{0pt}{1.1em}{0.3em}

\hypersetup{pdftitle={Turns}, pdfauthor={Robin Langer}, pdfborder={0 0 0}}

\title{\textbf{Turns}\\[2pt]\large What each kind of agent does in a month, and once a year}
\author{Robin Langer}
\date{7 November 2026}

\begin{document}
\maketitle

The economy advances one month at a time. Every month runs the \emph{monthly
turn}. Every twelfth month additionally runs the \emph{annual turn} first, and
then its monthly turn as usual. Each kind of agent layers its own actions onto a
common skeleton; this note takes them one at a time. A \textit{status} line flags
where the running metropolis loop does less than the design describes.

\section{The monthly turn --- a regular citizen}

A citizen who owns neither a firm nor a mine. Within one month, in order:

\begin{enumerate}
\item \textbf{Income.} If employed, receives his wage when his employer settles
its bill at the top of the month. A saver or idle citizen receives nothing here.
\item \textbf{Food.} Buys one unit of staple from the cheapest shop with stock.
If none can be had, he starves --- and may die.
\item \textbf{Wants.} Spends down his objective's chain --- nightclub
(happiness); luxury food then a cure (health); luxury food then tuition
(education) --- cheapest lot first, up to a budget of a quarter of his cash. A
\emph{savings}-objective citizen buys nothing here.
\item \textbf{Wellbeing.} This month's consumption becomes happiness or health;
last month's level decays.
\end{enumerate}

That is the whole turn for most of the town. Note the cash cascade: income
arrives first, he spends from it through the markets, and what he does not spend
stays as savings.

\section{The entrepreneur's extra turn}

An owner-operator of a producing firm (a shop, a luxury shop, a nightclub). He is
also a citizen and does everything in \S1. His \emph{firm} adds, across the same
month:

\begin{itemize}
\item \textbf{Opening (first).} The firm pays its bill: wages to its workers ---
himself included, which is where he collects his own wage as a citizen ---
machine maintenance, tax, and debt service (bank and private). If it cannot pay,
it folds. \emph{This is why income appears first in \S1: it is the firms paying
out.}
\item \textbf{Production.} The firm refills its shelf from its machine and its
crew, prices by the discovery rule (raise on a sellout, cut on a leftover, capped
by a cheaper rival), and posts the goods.
\item \textbf{Sales.} Citizens buy from the shelf during the month's markets; the
takings accrue to the firm.
\item \textbf{Remainder (last).} The \emph{draw}: the firm's cash above its
operating reserve is swept to him. This is his income as an owner, and it lands
at the very end of the month, after the markets have cleared.
\end{itemize}

So wages leave the firm first and the owner's draw is taken last; in between, the
households spend the wages back into the shops. His \emph{annual} extra actions
(investing in a machine, mortgaging his house) belong to the annual turn, \S7.

\section{The miner's extra turn}

A citizen who owns a mine --- a source of ore. Besides the citizen turn:

\begin{itemize}
\item \textbf{Sells ore} to a tech firm that needs raw material to build a
machine. The buyer pays the miner directly --- this is the miner's income --- and
the mine's stock falls by what it sold. The money reaching the owner (rather than
a bank) is the whole point: it is a producing firm's machine spend, passed through
the tech firm, landing in a mining household to be spent on luxury and
entertainment.
\item \textbf{Or, once a year, deposits the mine} at the bank for shares (a claim
on dividends). But a deposited mine belongs to the bank and can no longer sell
ore. So the miner chooses between a raw-material \emph{income flow} and an
\emph{equity claim}.
\end{itemize}

\section{The tech firm's extra turn}

A tech firm makes the machines that producing firms expand with. Its activity has
three parts: \textbf{buy} raw material (ore) from miners; \textbf{build} machines,
consuming that ore; and \textbf{sell} finished machines to producing firms. The
sale happens only at the year-end, because capacity is added once a year.

\subsection{Monthly or annual raw material? The coherent answer}

The awkward middle --- buy ore every \emph{month} but build only at \emph{year-end}
--- is indeed wrong: it would stockpile ore for a year with nothing built to show
for it. Two coherent designs:

\begin{itemize}
\item \textbf{(Recommended) Build through the year, sell at year-end.} Tech buys
ore \emph{as it builds}, month by month, and holds the finished machines as
inventory until the year-end sale. Ore purchase \,=\, building \,=\, monthly and
smooth; the miner earns a steady monthly income; only the \emph{sale} is annual.
\item \textbf{(Current simplification) Everything at year-end.} Tech buys the ore,
builds, and sells in a single step at the year-end. The miner's ore income is then
one annual lump.
\end{itemize}

The recommended design is both more realistic and better for demand: a steady
monthly income to the indigenous is spent steadily, rather than arriving in a lump
that then drains away. The running loop currently does the simplification.

\section{The bank's turn}

\begin{itemize}
\item \textbf{Monthly.} The bank earns: when a borrower settles the debt line of
his bill, the interest he pays raises the bank's equity (principal repayment does
not --- it just retires the loan). It also keeps its posted loan rate moving by
the scripted rule (a step toward the other bank's rate, a step up when stressed).
\item \textbf{Annually.} On the special month the bank \emph{lends} --- an
entrepreneur mortgages his house and the bank advances the principal; \emph{mints
shares} when a miner deposits a mine; and \emph{pays a dividend}, distributing its
standing equity to shareholders in proportion to their shares.
\end{itemize}

\noindent\textit{Status.} The bank's genesis equity is minus the whole money
supply (a pure deposit ledger with no reserve asset), so until that balance sheet
is changed, interest income cannot lift equity positive and the dividend never
actually pays. See \texttt{before-the-bank.pdf}.

\section{The government's turn}

\begin{itemize}
\item \textbf{Monthly.} Collects tax (each firm remits its tax in the opening
bill) and meets the payroll of the public firms it owns (the public universities
and hospitals), whose wages flow out in those firms' opening bills. Through its
public staff it supplies services --- doctors pointed at the diseases people hold,
teachers at the careers with listed slots.
\item \textbf{Annually.} May issue government debt in the capital market ---
borrowing from savers against future tax --- as a channel of its own.
\end{itemize}

\noindent\textit{Status.} In the running metropolis loop the government is
passive: tax is collected and public payroll is met through the opening bill, but
the active public-service assignment and government debt, though present in the
engine, are not yet driven here.

\section{The annual turn, and how it interleaves}

Every twelfth month is special. At the boundary into that month --- after the
previous month has fully settled, before the special month's own monthly turn
begins --- the annual turn runs, in this order:

\begin{enumerate}
\item \textbf{Bank dividend.} The bank distributes its standing earnings to
shareholders, so a household has income in hand before the next step.
\item \textbf{Children.} Households that can afford it bear a child; having the
dividend first suspends the need to liquidate assets to fund one.
\item \textbf{Capital market.} Equity (miners deposit mines for shares;
entrepreneurs mortgage houses and pay interest), then private debt (savers lend to
cash-short firms), then government debt.
\item \textbf{Real estate.} Houses change hands.
\item \textbf{Jobs and machines.} Tech sells its machines; a producing firm that
has been selling out, clears the margin, and can staff a machine buys one and
hires its crew.
\item \textbf{Rental.} Dwellings are let.
\end{enumerate}

Then the ordinary monthly turn for that month proceeds exactly as in \S1--\S4,
starting again with the opening bill and food.

\medskip
\noindent So the shape of time is:
\begin{center}
\begin{tabular}{@{}ll@{}}
\toprule
\textbf{An ordinary month} & monthly turn \\
\textbf{A year-boundary month} & \textbf{annual turn}, then monthly turn \\
\bottomrule
\end{tabular}
\end{center}

\noindent The annual turn reads the world as it stood at the close of the previous
month and sets up the capacity, demographics, and capital positions that the
special month --- and the eleven months after it --- then play out. Under the
recommended tech design (\S4.1) the ore-consuming \emph{building} happens every
month, while only the machine \emph{sale} sits in step~5 of the annual turn;
under the current simplification, building and the ore purchase sit there too.

\section{Implementation status at a glance}

Driven in the running loop: the monthly turn (\S1), the firm's opening bill,
production, and draw (\S2), invest day and --- behind flags --- private lending,
mine deposits, mortgages, interest, and the dividend attempt (\S7, steps 3 and 5).

Present in the engine but not yet driven in this loop: the real-estate market
(\S7.4), the rental market (\S7.6), government debt (\S6, \S7.3), the mining-land
trade, and monthly tech building (\S4.1). The dividend is driven but
accounting-blocked (\S5).

\end{document}
